Figure~\ref{fig:timeline} gives the overall picture before the
phase-by-phase results below: Iran's country-level network responsiveness,
measured continuously across the full study period, relative to its own
normal level. It is drawn in two panels because Baseline alone
(\PZeroStart{} to \PZeroBoundary{}) covers a period wider than the other
four phases put together, and within those four, Onset lasts only hours
-- too narrow to label on the same axis as Baseline without it and Prelude
overlapping. The top panel shows the full period for context, with the
post-baseline window shaded; the bottom panel re-plots exactly that shaded
window at its own scale, where all four remaining phases -- including the
hours-long Onset, marked as a highlighted band -- get legible, non-
overlapping labels. The five phases from Table~\ref{tab:phases} are marked
in both panels, along with the responsiveness level (\probingDarkRatio{}
of normal) below which a network counts as ``shut down'' in this paper's
terms. That specific level is not an arbitrary round number: it is the
most permissive cutoff that still almost never misfires on ordinary,
quiet-network noise, measured directly on networks in three countries with
no shutdown to detect, and checked against the entire study period rather
than just one sample month (an earlier, narrower version of this check was
tried first and rejected for exactly this reason --
Section~\ref{sec:state-derivation} has the full calibration). Three drops
are visible at a glance: a brief one in June 2025 (the reference event
inside the baseline phase), the previously unidentified 18-day event at
the start of the Prelude phase, and the sustained collapse that begins at
Onset and does not meaningfully recover until Restoration.

\begin{figure}[t]
  \centering
  \includegraphics[width=\linewidth]{figures/timeline}
  \caption{Country-level network responsiveness, relative to normal.
  \emph{Top:} full study period, with the post-baseline window (shaded)
  shown zoomed in below. \emph{Bottom:} Prelude through Restoration at
  higher time resolution; Onset (highlighted band) lasts only hours and
  would not be visible at the top panel's scale. In both panels, the
  dashed line marks the calibrated threshold below which a network counts
  as unreachable in this paper's analysis (\probingDarkRatio{} of its
  normal level; see Section~\ref{sec:state-derivation} for how this value
  was chosen); dotted vertical lines mark the phase boundaries from
  Table~\ref{tab:phases}.}
  \label{fig:timeline}
\end{figure}

\subsection{H1: mechanism (filtering vs.\ control-plane withdrawal)}
\label{sec:h1}

At every phase after baseline, most of the loss in reachability comes from
networks whose routes stay visible but that stop responding --
\emph{announced but dark} in this paper's terms -- not from routes
disappearing from the global routing table entirely (Figure~\ref{fig:state_share}).
The share of networks in this ``announced but dark'' state rises from
\darkSharePZero{} at baseline to \darkSharePOne{} during the Prelude phase
(driven by the January event discussed below, not a phase-wide effect), to
\darkSharePTwo{} at the shutdown's onset, and peaks at \darkSharePThree{}
during the sustained blackout, before falling to \darkSharePFour{} during
recovery. The share of networks actually withdrawn from routing, by
contrast, moves far less across the same span
(\withdrawnSharePZero{} $\to$ \withdrawnSharePOne{} $\to$ \withdrawnSharePTwo{}
$\to$ \withdrawnSharePThree{} $\to$ \withdrawnSharePFour{}). This pattern --
routes staying announced while reachability collapses -- is
\textbf{observed} directly in the measurement.

\begin{figure}[t]
  \centering
  \includegraphics[width=0.85\linewidth]{figures/state_share}
  \caption{Share of networks in each connectivity state, by phase.
  ``Announced but dark'' rises and falls with the phases; ``withdrawn'' stays
  comparatively flat throughout.}
  \label{fig:state_share}
\end{figure}

Two independent checks support treating this as a real, phase-linked
effect rather than a measurement artifact. First, the same measurement
applied to the \nControlAsnsTotal{}-network control population (three
countries with no documented shutdown) does not track Iran's phase
structure at all: the control group's share of ``dark'' networks stays low
throughout (average \controlMeanDarkShare{}, never above
\controlMaxDarkShare{}, across \nControlBins{} time points checked), rising
only slightly and only during the sustained blackout and recovery phases,
never approaching Iran's peak (Figure~\ref{fig:control_falsification}).
Only \nControlArtifactBins{} of those \nControlBins{} time points cross the
threshold set in advance for flagging a possible artifact, and all of them
fall late in the study period, unrelated to Iran's own rise -- this is
\textbf{observed} to not explain Iran's pattern. Second, the exact timing
of the shutdown's onset can be checked two ways at once. The
country-level responsiveness signal crosses into ``dark'' territory at
\PTwoStart{} (Section~\ref{sec:phases}), while the bulk of routing
withdrawals -- the point by which 5\% of affected address blocks had
already been withdrawn, counting only withdrawals from that moment onward
so ordinary background routing changes beforehand are not mistaken for
part of the event -- does not happen until \withdrawalWaveTPFive{}, roughly
\withdrawalWaveGapMinutes{} minutes later (\withdrawalWaveNPrefixes{}
address blocks withdrawn in that wave). That responsiveness collapses
before the bulk of routing withdrawals is \textbf{observed} directly;
that this reflects traffic filtering happening ahead of a separate, slower
routing-level action is \textbf{strongly implied} by two independent
measurements (routing state and responsiveness) agreeing on the same
sequence, with no plausible artifact explanation surviving the control
check above. This reading is further corroborated by an independently
authored study using an entirely different measurement approach (passive
network scans and active connection probing from external vantage points,
with no overlap with this paper's own data sources), which reports the
2026 blackout enforced through traffic-blocking at a small number of
central points while routing announcements themselves stayed
stable across 96.5--97.4\% of the Iranian address space it
checked~\citep{jahromi2026multiperspective} -- a more mechanistically
specific description of the same routes-stay-visible pattern this paper
measures.

\begin{figure}[t]
  \centering
  \includegraphics[width=0.7\linewidth]{figures/control_falsification}
  \caption{Mean share of ``dark'' networks, Iran vs.\ the control population, by phase. The
  control group does not track Iran's rise during the shutdown.}
  \label{fig:control_falsification}
\end{figure}

A closer look at the individual measurements taken right around the
shutdown's onset (Figure~\ref{fig:transition_zoom}) shows a single sharp
step, with the level immediately after the step already matching the
level the blackout holds at for months afterward -- \textbf{observed}, and
the reason this paper treats the onset phase's reported numbers as
representative of a real, stable state rather than a fluke of exactly
when the measurement happened to be taken (see Section~\ref{sec:phases}
for why this phase is unusually short).

\begin{figure}[t]
  \centering
  \includegraphics[width=0.7\linewidth]{figures/transition_zoom}
  \caption{Share of networks ``dark'' and ``withdrawn'' at four consecutive
  measurements spanning the shutdown's onset.}
  \label{fig:transition_zoom}
\end{figure}

\paragraph{The January 2026 event.} Scanning the full study period for
sustained drops in connectivity (Section~\ref{sec:phases}) surfaces an
18-day near-total collapse beginning \POneStart{}, entirely inside what a
press-reported timeline would treat as one undifferentiated ``throttling''
period, with no detectable change at all around the commonly reported
late-December-2025 start date for that period, in any signal checked. This
event is \textbf{observed} in this study's own measurement and passes the
same control-population check as the main shutdown. It is not, however, a
novel discovery: it was one of the most heavily publicly reported
shutdowns of the period, and a separate, contemporaneous measurement study
covering the same window, using an overlapping combination of routing data
and outage-monitoring methods, already exists~\citep{aceto2026iran}. This
paper's contribution regarding that specific event is an independent
reproduction via measurement, with finer routing-level timing than prior
public coverage offered, not a claim of discovery.

\subsection{H2: topology of disconnection}
\label{sec:h2}

How often a network's set of upstream connections changes -- a proxy for
reconfiguration activity in the network's underlying wiring -- spikes at
the shutdown's onset (\transitionRatePTwo{} of networks changing upstreams
in that period, several times the \transitionRatePZero{} baseline rate) and
then drops to the lowest rate of the whole study during the sustained
blackout (\transitionRatePThree{}), with a further small decline during
recovery (\transitionRatePFour{}) (Figure~\ref{fig:h2_churn}). This
spike-then-flat pattern is \textbf{observed}: reconfiguration activity is
concentrated at the moment of the flip into blackout, not sustained
throughout it, consistent with a one-time change rather than continued
active rerouting during the months the blackout lasted.

\begin{figure}[t]
  \centering
  \includegraphics[width=0.7\linewidth]{figures/h2_churn}
  \caption{Rate of upstream-connection changes, by phase.}
  \label{fig:h2_churn}
\end{figure}

This measurement does not, on its own, establish \emph{where in the
network} this reconfiguration happened. The original question -- whether
disconnection was carried out at a small number of international gateway
networks, rather than distributed across many individual access providers
-- would require breaking this measurement down by each network's role,
which this analysis has not yet done at scale. A small, illustrative check
on the two networks known to carry Iran's international gateway traffic
found that one changed its upstream connections during the onset window and
one did not, out of only four total measurements -- far too small a sample
to support a conclusion either way, and reported here only so that this
gap is explicit rather than quietly absorbed into the aggregate finding
above. Resolving this specific question is left as an open limitation
(Section~\ref{sec:discussion}), not folded into the aggregate result.

\subsection{H3: selectivity of restoration}
\label{sec:h3}

Routing visibility, broken down by network category, stays at or above
92\% for every category in every phase, including during the blackout
itself (Figure~\ref{fig:h3_visibility_by_type}) -- \textbf{observed}, and
consistent with the filtering-not-withdrawal reading from H1: routes
mostly stay visible across every sector; sites simply stop responding when
checked. The mobile category sits persistently about 4--5 percentage
points below every other category, including during ordinary
pre-shutdown months, not specifically during the blackout --
\textbf{observed} as a structural, year-round difference rather than a
blackout-specific effect.

\begin{figure}[t]
  \centering
  \includegraphics[width=0.8\linewidth]{figures/h3_visibility_by_type}
  \caption{Mean routing visibility by network category, monthly. All other
  categories (shaded band) stay within a narrow range; mobile sits
  persistently lower and the unclassified group is the most volatile.}
  \label{fig:h3_visibility_by_type}
\end{figure}

The centerpiece result for restoration speed (Figure~\ref{fig:h3_restoration_speed})
uses the fine-grained, minute-level event data described in
Section~\ref{sec:methods} rather than the routine snapshot measurements
(which are too coarse to resolve timing within the recovery window -- most
address blocks appear ``restored'' at the very first snapshot taken after
recovery began, an artifact of the measurement's resolution rather than
evidence that recovery was instantaneous). By category, the typical delay
between the start of the recovery phase and an individual network coming
back online ranges from \restoreFastestMin{} minutes
(\restoreFastestType{}) to \restoreSlowestMin{} minutes
(\restoreSlowestType{}), with the state telecom operator at
\restoreDelayMinStateTelecom{} minutes and a middle tier -- education,
media, internet providers, and hosting -- clustered around 13--15 minutes.
This ordering is \textbf{observed}.

\begin{figure}[t]
  \centering
  \includegraphics[width=0.7\linewidth]{figures/h3_restoration_speed}
  \caption{Typical delay before a network came back online after the start
  of the recovery phase, by category (fine-grained, minute-level
  measurement).}
  \label{fig:h3_restoration_speed}
\end{figure}

This is a \textbf{mixed} result relative to the hypothesis that recovery
was prioritized by sector, not a clean confirmation, and is reported as
such. In its favor: government networks restore fastest and most
consistently, and the unclassified, long-tail group of networks is both
slowest to restore and least likely to restore at all
(\neverRestoredUnclassified{} of unclassified address blocks never come
back online by the end of the study period, versus \neverRestoredGovernment{}
for government and \neverRestoredStateTelecom{} for the state telecom
operator). Complicating a simple priority story: the mobile category has a
fast typical delay but a long tail of slow stragglers. A related
finding -- that certain categories of network are effectively protected
from disruption -- is independently corroborated: the external study cited
above separately reports that Iranian academic networks and the country's
largest commercial content-delivery network retained substantially higher
visibility than other Iranian networks throughout the
shutdown~\citep{jahromi2026multiperspective}, consistent with (though not
identical in scope to) this paper's own always-reachable group of hosting
and content-delivery networks. Classification quality is externally
validated (Cohen's $\kappa$ = \kappaValue{}, \kappaAgreePct{} agreement on
a blind sample of \kappaN{} networks), but the single correction above is
a reminder that category-level results are only as reliable as the
underlying categorization, not a reason to distrust the overall pattern.

\subsection{H4: comparison to prior shutdown/restoration events}
\label{sec:h4}

Five historical events -- the November 2019 shutdown, the June 2025
shutdown, the January 2026 collapse (Section~\ref{sec:h1}), the February
2026 onset, and the May 2026 recovery -- can be compared for how quickly
they unfolded, using fine-grained routing data for each
(Figure~\ref{fig:h4_comparison}). The primary comparison uses the range
between the 5th and 95th percentile of when individual address blocks
were first withdrawn from routing -- a measure that is not thrown off by
a single unusually early or late outlier. By this measure, the February
2026 onset was the fastest of the four collapse-type events, completing in
\durationFebOnsetRangeHours{} hours, against \durationNovRangeHours{} hours
for November 2019, \durationJunRangeHours{} hours for June 2025, and
\durationJanEventRangeHours{} hours -- by far the widest of the four -- for
the January 2026 collapse. That last figure is \textbf{observed} and fits
the pattern already established in Section~\ref{sec:h1}: the January event
was identified there as an 18-day near-total collapse rather than a sharp
single-day cutover, and its withdrawal timing here is consistent with that
same gradual, rolling character rather than a discrete event like the
February onset.

The typical (median) delay tells a different story, however, and this
time it is the January 2026 collapse that comes out fastest
(\durationJanEventMedianHours{} hours), narrowly ahead of June 2025
(\durationJunMedianHours{} hours) and well ahead of November 2019
(\durationNovMedianHours{} hours) and the February 2026 onset
(\durationFebOnsetMedianHours{} hours, the slowest of the four by this
measure). So the February 2026 onset is simultaneously the fastest event
by overall range and the slowest by typical delay -- and the January 2026
collapse is the reverse, slowest by range and fastest by typical delay.
Both facts are \textbf{observed}; neither is smoothed into a single
``which event was faster'' headline, since the two measures capture
different things (how tightly clustered the bulk of withdrawals was,
versus how long the median network waited) and rank the events
differently depending on which one is asked.

\begin{figure}[t]
  \centering
  \includegraphics[width=0.8\linewidth]{figures/h4_comparison}
  \caption{Speed of each event: 5th-95th-percentile range vs.\ typical
  (median) delay.}
  \label{fig:h4_comparison}
\end{figure}

The May 2026 recovery window shows about as many withdrawals as the onset
windows do (\nWithdrawnMayRestoration{} address blocks) despite being a
\emph{recovery} event -- \textbf{observed}, and consistent with (though
not independently confirmed beyond) the uneven, partial character of the
recovery phase described in Section~\ref{sec:phases}: a mix of networks
coming online and going back offline during partial recovery, rather than
one clean, one-directional wave of restoration.

This comparison has no equivalent to the control-population check used
elsewhere in this paper: it compares Iran to itself across seven years of
history rather than to other countries at a fixed point in time, so
changes in the measurement infrastructure itself over that span cannot be
ruled out the way they can for the within-study-period results above. This
is reported as an open limitation, not an implicit claim that all five
events are perfectly comparable.
